\documentclass[11pt]{amsart}
\usepackage[localtheorems]{raeez-math-template}

\providecommand{\bC}{\mathbb{C}}
\providecommand{\bP}{\mathbb{P}}
\providecommand{\bZ}{\mathbb{Z}}
\providecommand{\cA}{\mathcal{A}}
\providecommand{\cO}{\mathcal{O}}
\providecommand{\CoHA}{\mathrm{CoHA}}
\providecommand{\Tot}{\operatorname{Tot}}
\providecommand{\PhiFA}{\Phi^{\mathrm{FA}}}
\providecommand{\Perf}{\mathrm{Perf}}
\providecommand{\CY}{\mathrm{CY}}

\newtheorem{theorem}{Theorem}
\newtheorem{proposition}[theorem]{Proposition}
\newtheorem{lemma}[theorem]{Lemma}
\newtheorem{corollary}[theorem]{Corollary}
\newtheorem{conjecture}[theorem]{Conjecture}
\theoremstyle{definition}
\newtheorem{definition}[theorem]{Definition}
\theoremstyle{remark}
\newtheorem{remark}[theorem]{Remark}

\title[Compact Hodge trace stratification]{Compact Hodge trace
stratification in the Calabi--Yau examples}
\author{Raeez Lorgat}
\date{2026-04-30}

\begin{document}

\begin{abstract}
The compact Hodge supertrace
\[
        \Xi(X)=\sum_q(-1)^q h^{0,q}(X)
\]
is the compact numerical invariant of the Calabi--Yau examples. Its
interpretation as \(\kappa_{\mathrm{ch}}\) is a theorem on the constructed
\(\Phi_d\)-loci and a construction problem beyond them. The table through
dimension five separates this compact invariant from the noncompact local
surface value, the conifold NCCR value, the Borcherds weight, and the fibre
Euler characteristic. These numbers come from different constructions and
must carry their subscripts.
\end{abstract}
\maketitle

\section{Four invariants}

\begin{definition}
For a compact Calabi--Yau \(X\) define
\[
        \Xi(X)=\sum_q(-1)^q h^{0,q}(X).
\]
On the compact loci where the trace comparison for
\(A_X=\Phi_d(\Perf(X))\) has been constructed, this gives
\(\kappa_{\mathrm{ch}}(A_X)=\Xi(X)\). The \(d\geq4\) interpretation remains
conditional. Heisenberg, BKM, fibre, BCOV, and noncompact toric
specialisations are separate invariants. The
categorical invariant is
\[
        \kappa_{\mathrm{cat}}(X)=\chi(\cO_X),
\]
with K\"unneth-multiplicative behaviour on products. The Borcherds invariant
is
\[
        \kappa_{\mathrm{BKM}}(\Phi_N)=\frac{c_N(0)}{2},
\]
on the CHL/Gritsenko--Nikulin Borcherds-product levels
\(N\in\{1,2,3,4,6\}\) with the chosen weak Jacobi input. The
Cl\'ery--Gritsenko diagonal-divisor catalogue is a different
Borcherds-product catalogue. In its triple order the weights are
\[
        \left(5,2,3,1,2,\frac12,\frac32,1\right),
\]
with integer-weight paramodular rows and two half-integer multiplier-system
rows. These weights are not the CHL sequence \((5,3,2,\tfrac32,1)\). The fibre
invariant \(\kappa_{\mathrm{fiber}}\) is the topological Euler/Mukai-lattice
rank of the chosen fibre; for a K3 fibre it is \(24\), not
\(\chi(\cO_{K3})=2\).
\end{definition}

\section{Compact dimension table}

\begin{theorem}
The Hodge supertrace gives the following unconditional compact table:
\[
\begin{array}{c|c|c|c}
d & X & (h^{0,0},\ldots,h^{0,d}) & \Xi(X)\\
\hline
1 & E & (1,1) & 0\\
2 & K3 & (1,0,1) & 2\\
2 & \text{abelian surface} & (1,2,1) & 0\\
3 & \text{quintic threefold} & (1,0,0,1) & 0\\
3 & K3\times E & (1,1,1,1) & 0\\
3 & E^3 & (1,3,3,1) & 0\\
4 & \text{sextic fourfold} & (1,0,0,0,1) & 2\\
5 & \text{generic strict }\CY_5 & (1,0,0,0,0,1) & 0
\end{array}
\]
\end{theorem}

\begin{proof}
For a strict compact Calabi--Yau \(d\)-fold, Serre duality and the standard
Hodge-column vanishing leave only \(h^{0,0}\) and \(h^{0,d}\), giving
\(1+(-1)^d\). Product examples follow by K\"unneth on \(h^{0,*}\); this gives
the \(K3\times E\) and \(E^3\) rows. The abelian-surface row is the exterior
algebra on \(H^1(\cO)\). The hypersurface rows use Lefschetz and Kodaira
vanishing. The displayed values are the resulting Hodge supertraces. The
constructed-locus theorem identifies these numbers with
\(\kappa_{\mathrm{ch}}\) only where the object \(A_X=\Phi_d(\Perf(X))\) and
its trace comparison have been built; the \(d\geq4\) rows remain compact
Hodge data until that construction is supplied.
\end{proof}

\section{Local and singular examples}

\begin{proposition}[Local \(\bP^2\)]
For the local surface \(X=\Tot(K_{\bP^2}\to \bP^2)\) the chiral shadow
normalisation gives
\[
        \kappa_{\mathrm{ch}}(X)=\frac32 .
\]
\end{proposition}

\begin{proof}
The local \(\bP^2\) shadow theorem computes the same value by the
\(\bZ/3\)-McKay quiver datum for \(\bC^3/\bZ_3\), by three affine toric
charts each contributing \(1/2\), by the MacMahon exponent
\(M(q)^{\chi_{\mathrm{top}}(\bP^2)}\), and by the compactly supported
Dolbeault trace on the base. The trace comparison carries the local
compact-support normalisation, which contributes the factor \(1/2\).
Thus the local-surface value is
\(\chi_{\mathrm{top}}(\bP^2)/2=3/2\). This is a noncompact \(\CY_3\)
local-surface reading, separate from the compact Hodge supertrace table.
\end{proof}

\begin{proposition}[Conifold]
The resolved conifold
\[
        \Tot(\cO(-1)\oplus\cO(-1)\to\bP^1)
\]
is outside the local surface family \(\Tot(K_S\to S)\). The conifold NCCR
calculation gives
\[
        \kappa_{\mathrm{ch}}(\mathrm{conifold})=1 .
\]
\end{proposition}

\begin{proof}
The resolved conifold is the small resolution over \(\bP^1\), outside the
canonical-bundle-over-a-surface family. The finite-quotient BKR McKay route
for \(\bC^3/G\) is inapplicable: the conifold is the affine hypersurface
\(xy-zw=0\), and its crepant resolution is governed by the
Szendroi/Klebanov--Witten noncommutative crepant resolution, equivalently
the two-vertex conifold quiver with potential. The value \(1\) is the
direct conifold NCCR/Hochschild shadow calculation: the level-one
DT/CoHA coefficient on the compact class \([\bP^1]\) is \(1\). The
half-Euler value of the exceptional curve is only a numerical coincidence
here.
\end{proof}

\begin{proposition}[Affine positive half and toric gluing]
For the affine chart \(\bC^3\), with the three-loop Jordan quiver and
potential \(W=\operatorname{tr}(XYZ-XZY)\),
\[
        \CoHA^{\mathrm{sph}}_{\mathrm{crit},T}(\bC^3,W)
        \cong
        Y^+(\widehat{\mathfrak{gl}}_1).
\]
The target is the positive half. It is not
\(\mathcal W_{1+\infty}\). A finite toric CY$_3$ atlas with
simplex-compatible DWR/Ran Rees hCS--Hall data glues these local positive
halves to
\[
        H^+_{X_\Sigma}
        =
        \operatorname{hocolim}_{\alpha\in\Sigma(3)}
        \CoHA^{\mathrm{sph}}_{\mathrm{crit},T}(U_\alpha).
\]
Compact critical CoHA, quasi-NCCR, and Hall--Drinfeld double outputs are
separate compact realisation problems: they require monoidal
vanishing-cycle realisation, compact support, a non-degenerate pairing
after radical quotient, completion, opposite and Cartan halves, and
PBW/no-extra-relations compatibility.
\end{proposition}

\begin{proof}
Schiffmann--Vasserot identify the spherical equivariant critical Hall
algebra of the tripled Jordan quiver with the shuffle presentation of
\(Y^+(\widehat{\mathfrak{gl}}_1)\). The DWR/Ran Rees comparison is formed
chartwise from finite cyclic critical atlases; Stokes' formula on the
simplex variables gives the face compatibility, and Hall wall-crossing
quasi-isomorphisms identify the overlaps. Taking the finite homotopy
colimit therefore gives the displayed positive half. None of these
operations constructs the continuous dual, Hopf pairing, compact-support
vanishing-cycle functor, or completed double needed for the compact
objects.
\end{proof}

\section{Borcherds calculation}

\begin{theorem}[Borcherds weight]
For the CHL/Gritsenko--Nikulin Borcherds-product levels
\(N\in\{1,2,3,4,6\}\) with weak Jacobi input \(\phi_N\),
\[
        \kappa_{\mathrm{BKM}}(\Phi_N)=\frac{c_N(0)}{2}.
\]
The constants are
\[
\begin{array}{c|ccccc}
N & 1 & 2 & 3 & 4 & 6\\
\hline
c_N(0) & 10 & 8 & 6 & 4 & 2\\
\kappa_{\mathrm{BKM}}(\Phi_N) & 5 & 4 & 3 & 2 & 1
\end{array}
\]
\end{theorem}

\begin{proof}
In Borcherds' singular theta-lift normalisation, the automorphic product
attached to a weight-zero weak Jacobi input has weight equal to one half of
the constant coefficient of that input. In the CHL/Gritsenko--Nikulin
levels this coefficient is \(c_N(0)\), giving the displayed table. At
\(N=1\) the input \(\phi_{0,1}\) has \(c_1(0)=10\), so the theta-normalised
Igusa square root has weight \(5\). The denominator identity reads that
weight as \(\kappa_{\mathrm{BKM}}\). This calculation belongs to the
automorphic Borcherds construction and is independent of the compact Hodge
supertrace. Cl\'ery--Gritsenko diagonal-divisor rows use a separate catalogue;
the comparison tuple \((5,2,1,1,\frac12,1,\frac14,0)\) is not that
catalogue, and no Cl\'ery--Gritsenko row has weight \(\frac14\) or \(0\).
\end{proof}

\begin{proposition}[Cl\'ery--Gritsenko diagonal-divisor rows]
The Cl\'ery--Gritsenko classification gives exactly eight diagonal-divisor
genus-two Borcherds products:
\[
\begin{array}{c|c|c|c}
(t,N) & F_{t,N} & \mathrm{wt}(F_{t,N}) & 2\,\mathrm{wt}(F_{t,N})\\
\hline
(1,1) & \Delta_5 & 5 & 10\\
(2,1) & \Delta_2 & 2 & 4\\
(1,2) & \nabla_3 & 3 & 6\\
(3,1) & \Delta_1 & 1 & 2\\
(1,3) & \nabla_2 & 2 & 4\\
(4,1) & \Delta_{1/2} & \frac12 & 1\\
(1,4) & \nabla_{3/2} & \frac32 & 3\\
(2,2) & Q_1 & 1 & 2
\end{array}
\]
For each row, the Borcherds invariant of the product is its weight:
\(\kappa_{\mathrm{BKM}}(F_{t,N})=\mathrm{wt}(F_{t,N})\). The last column is
the corresponding twice-weight tuple \((10,4,6,2,4,1,3,2)\), not the CHL
constant table \((10,6,4,3,2)\).
\end{proposition}

\begin{proof}
Cl\'ery--Gritsenko classify the reflective genus-two modular forms with
diagonal divisor and obtain exactly these eight holomorphic products. The
candidate \((2,4;\frac12)\) is not holomorphic, so it is not a ninth row.
Their Borcherds product theorem gives the product weight as one half of the
cusp-weighted sum of the row's zeroth Jacobi coefficients; only the
single-cusp cases reduce to a single coefficient \(c(0,0)/2\). The CHL
ladder instead uses the five weak Jacobi inputs \(\phi_N\),
\(N\in\{1,2,3,4,6\}\), producing Borcherds products \(\Phi_N\) with
constants \((10,6,4,3,2)\) and weights \((5,3,2,\tfrac32,1)\)
(Jatkar--Sen; Govindarajan--Krishna; the once-recorded ladder
\((5,4,3,2,1)\) is retracted). The two lists share
the \((1,1)\) product \(\Delta_5=\Phi_1\) directly. The scalar-square identities
that occur in the CHL/dyon literature are row-specific:
\[
        \nabla_3^2=\Phi_6^{(1,2)},\qquad
        \nabla_2^2=\Phi_4^{(1,3)},\qquad
        \nabla_{3/2}^2=\Phi_3,\qquad
        Q_1^2=\Phi_2^{(2,4)}.
\]
They use their own levels, covers, branches, and denominator data. They
match the remaining CHL square-root denominators with the diagonal-slice
rows \((1,2),(1,3),(1,4),(2,2)\) row by row; they do
not identify the five-entry CHL table with the full eight-row
Cl\'ery--Gritsenko catalogue,
and they do not turn the twice-weight tuple
\((10,4,6,2,4,1,3,2)\) into the CHL constant table \((10,6,4,3,2)\).
\end{proof}

\begin{proposition}[\(K3\times E\) spectrum]
The four labelled \(K3\times E\) values are
\[
\begin{aligned}
\kappa_{\mathrm{cat}}(K3\times E)&=0,\\
\kappa_{\mathrm{ch}}^{\mathrm{Heis}}(K3\times E)&=3,\\
\kappa_{\mathrm{BKM}}(\Phi_1)&=5,\\
\kappa_{\mathrm{fiber}}(K3)&=24.
\end{aligned}
\]
Thus the labelled spectrum is \(\{0,3,5,24\}\). The labels are part of the
assertion.
\end{proposition}

\begin{proof}
K\"unneth gives
\[
        \Xi(K3\times E)=(1-0+1)(1-1)=0
\]
and
\[
        \kappa_{\mathrm{cat}}(K3\times E)
        =\chi(\cO_{K3})\chi(\cO_E)=2\cdot 0=0.
\]
The Heisenberg--Mukai specialisation is additive on the K3 and elliptic
branches, giving
\[
        \kappa_{\mathrm{ch}}^{\mathrm{Heis}}(K3\times E)
        =\kappa_{\mathrm{ch}}^{\mathrm{Heis}}(K3)
        +\kappa_{\mathrm{ch}}^{\mathrm{Heis}}(E)
        =2+1=3.
\]
The \(N=1\) Borcherds row gives
\(\kappa_{\mathrm{BKM}}(\Phi_1)=c_1(0)/2=10/2=5\). The fibre invariant is
\[
        \kappa_{\mathrm{fiber}}(K3)=\operatorname{rk}H^*(K3,\bZ)=24.
\]
The number \(2\) is the fibre categorical value \(\chi(\cO_{K3})\), not the
total-space value \(\kappa_{\mathrm{cat}}(K3\times E)\). The equality
\(3+2=5\) is a coincidence in this row, not a formula for
\(\kappa_{\mathrm{BKM}}\).
\end{proof}

\begin{proposition}[Igusa square normalisations]
Let \(\Delta_5\) denote the theta-constant normalisation of the Igusa
square root, so that
\[
        [q^{1/2}r^{1/2}s^{1/2}]\Delta_5=64 .
\]
Let \(D_5\) be the monic Borcherds product attached to the
\(\phi_{0,1}\) input. Then
\[
        D_5=64^{-1}\Delta_5 .
\]
Oberdieck--Pixton's primitive scalar product squares this monic product:
\[
        \Phi_{10}^{\mathrm{OP}}
        =D_5^2
        =4096^{-1}\Delta_5^2.
\]
In the quotient-by-\(E\), Behrend-weighted \(K3\times E\) reduced
DT/PT normalisation,
\[
        Z^X_{\mathrm{OP}}
        =Z^X_{\mathrm{DT}}
        =-\bigl(\Phi_{10}^{\mathrm{OP}}\bigr)^{-1}
        =-4096\,\Delta_5^{-2}.
\]
The DVV/DMVV physics normalisation uses the unscaled Igusa square up to
scalar; in the theta-normalised convention this is
\[
        \Phi_{10}^{\mathrm{un}}=\Delta_5^2 .
\]
\end{proposition}

\begin{proof}
The Borcherds--Gritsenko--Nikulin product supplies the monic product
\[
        D_5
        =
        q^{1/2}r^{1/2}s^{1/2}
        \prod(1-q^nr^ls^m)^{f(nm,l)}
\]
with vacuum coefficient \(1\), before expanding the real-root factor. The
theta-constant normalisation of the same product has coefficient \(64\) at
\(q^{1/2}r^{1/2}s^{1/2}\). Hence \(D_5=64^{-1}\Delta_5\). Squaring gives
\[
        D_5^2=(64^{-1}\Delta_5)^2=4096^{-1}\Delta_5^2 .
\]
Oberdieck--Pixton's primitive scalar theorem gives the negative inverse of
this monic square. The reduced DT equality uses the quotient-by-\(E\),
Behrend-weighted \(K3\times E\) comparison through reduced DT/PT and
stable-pairs/Gromov--Witten correspondence. In that normalisation,
\[
        Z^X_{\mathrm{OP}}=Z^X_{\mathrm{DT}}
        =-(D_5^2)^{-1}=-4096\,\Delta_5^{-2}.
\]
DVV/DMVV use the genus-two Igusa form up to scalar; the unscaled
theta-normalised automorphic square is
\(\Phi_{10}^{\mathrm{un}}=\Delta_5^2\).
The scalar \(4096=64^2\) is a normalisation conversion, not a new
Borcherds weight and not a Hall-orientation character.
\end{proof}

\begin{remark}
The compact Hodge trace, the Heisenberg--Mukai specialisation, the
Borcherds weight, and the fibre rank are four constructions. The notation
records this separation; erasing the subscripts erases the theorem.
\end{remark}

\end{document}
